\chapter{Apresentação da Ferramenta}

\section{O que é o UnDF-TeX}

O UnDF-TeX é uma classe e um conjunto de arquivos LaTeX destinados à preparação de Trabalhos de Conclusão de Curso. O projeto separa conteúdo, metadados, elementos opcionais, bibliografia e implementação interna para reduzir formatação manual e tornar a saída reproduzível.

O template gera capa, folha de rosto, folha de aprovação, resumos, sumário, capítulos, referências e elementos opcionais a partir dos dados e arquivos preenchidos pelo usuário. Também oferece comandos para citações, figuras, tabelas, fluxogramas, equações, pseudocódigos e referências cruzadas.

\section{Finalidade e limites}

O objetivo é automatizar tarefas repetitivas e detectar problemas técnicos comuns. O template não avalia mérito acadêmico, não verifica a veracidade das fontes, não decide se uma pesquisa precisa de autorização ética e não substitui revisão humana.

\manualnote{Automação de aparência não transforma uma escolha do projeto em regra da ABNT nem em decisão institucional. O manual acadêmico e as orientações do curso continuam sendo as referências para conteúdo e aplicabilidade.}

\section{Versões e situação}

Este guia corresponde à revisão documental \texttt{1.0.0} e descreve o UnDF-TeX \texttt{1.0.0}. Antes de iniciar ou atualizar um trabalho, confira o arquivo \texttt{VERSION.txt}, o \href{https://github.com/igormrqs/undftex-publico/blob/main/CHANGELOG.md}{histórico de alterações} e a origem do pacote recebido. A versão do manual e a versão técnica do template têm controles próprios.

Versões do projeto podem mudar. Preserve uma cópia funcional do trabalho antes de substituir a classe ou os módulos internos e faça nova regressão do PDF depois da atualização.

\section{Relação com o manual acadêmico}

Este guia remete ao Manual de TCC sempre que o usuário precisa decidir se um elemento é obrigatório, opcional ou adequado. Aqui são documentados apenas localização, sintaxe, processamento, mensagens e comportamento técnico.
